Este trabalho investiga se a geração de uma representação intermediária comprimida, seguida de expansão determinística, produz modelos BPMN 2.0 válidos de forma mais confiável e mais econômica em \textit{tokens} do que a geração direta de XML por modelos de linguagem. Projetou-se uma linguagem de domínio específico com gramática EBNF formalmente especificada e um transpilador determinístico que a expande em XML validado contra o esquema oficial da OMG. Construiu-se um corpus de 1.021 processos a partir de fontes públicas de descrições textuais e especializou-se um modelo de sete bilhões de parâmetros por ajuste supervisionado com adaptação de baixo posto sobre pesos quantizados. A avaliação seguiu protocolo congelado antes de qualquer execução, com sete braços experimentais, 53 processos de referência modelados por especialistas e três repetições por item, totalizando 1.113 gerações. O componente determinístico produziu XML válido em 100\% das amostras e preservou a equivalência topológica em 99,6\%, e a representação proposta mostrou-se seis vezes mais compacta que o XML lógico equivalente. Contrariando a hipótese principal, fornecer a gramática no \textit{prompt} reduziu a confiabilidade: modelos de fronteira assim instruídos produziram saída válida em 64,8\% e 56,6\% das gerações, contra 98,1\% e 94,3\% da geração direta de XML pelos mesmos modelos, com significância estatística e replicação em duas famílias distintas. A configuração que este trabalho propõe — um modelo pequeno com a notação internalizada por ajuste fino, e não descrita no \textit{prompt} — atingiu 86,8\% de validade, vinte e dois pontos percentuais acima do modelo de fronteira instruído pela mesma gramática, ainda que abaixo dos 98,1\% da geração direta de XML, emitindo 200 \textit{tokens} medianos contra 631, a um custo de treinamento inferior a um dólar; patamar próximo, de 83,0\%, replicou-se em verificação exploratória sobre duas famílias adicionais de modelos. Em fidelidade topológica, porém, o modelo especializado empata estatisticamente com o modelo de fronteira instruído por \textit{prompt} e não alcança a geração direta de XML. Mediu-se ainda a concordância entre modelos de referência produzidos por especialistas distintos para um mesmo processo, obtendo-se F1 estrutural de 0,1449, patamar em que todos os métodos avaliados operam; disso decorre que métricas estruturais que identificam nós por rótulo saturam próximo ao ruído e não devem ser lidas como percentual de acerto. Conclui-se que o benefício de representações intermediárias compactas não se realiza por instrução em contexto, e sim por especialização do modelo, o que torna o ajuste fino condição de viabilidade da abordagem, e não etapa acessória.

\textbf{Palavras-chave}: BPMN; linguagem de domínio específico; modelos de linguagem; ajuste fino supervisionado; compressão de \textit{tokens}.
